\documentclass[a4paper,12pt]{article}
\usepackage[utf8]{inputenc}
\usepackage[T1]{fontenc}
\usepackage{lmodern}
\usepackage{geometry}
\geometry{
 a4paper,
 left=1.25in,
 right=1.25in,
 top=1in,
 }
\usepackage[croatian,english]{babel}    %za hrvatske naslove
\usepackage[nottoc]{tocbibind}  %za pravilan table of content
\usepackage{graphicx}   %za dodavanje slika
\usepackage{amsmath}    %za matematičke forumele
\usepackage{subcaption} %za dvije slike u jednoj
\usepackage{booktabs}   %za bolje tablice
\usepackage{multirow}   %za multirow tablice
\usepackage{cite}
\usepackage{url}       %za URL linkove u bibliografiji
\usepackage{float}
\usepackage{hyperref}  %za hiperlinkove - mora biti zadnji
\usepackage{listings}
\usepackage{xcolor}
\usepackage{listingsutf8}

% Configure listings for Python
\lstset{
    language=Python,
    basicstyle=\ttfamily\small,
    keywordstyle=\color{cyan},
    commentstyle=\color{green!70!white},
    stringstyle=\color{orange},
    numberstyle=\tiny\color{gray!70!white},
    showstringspaces=false,
    breaklines=true,
    frame=single,
    numbers=left,
    inputencoding=utf8,
    extendedchars=true
}

% Konfiguracija hyperref-a za boje i stil linkova
\hypersetup{
    colorlinks=true,        % linkovi su obojeni umjesto u okvirima
    linkcolor=black,         % boja unutarnjih linkova (TOC, reference)
    citecolor=black,        % boja citata
    urlcolor=blue,          % boja URL linkova
    filecolor=magenta,      % boja linkova na datoteke
    menucolor=red,          % boja linkova na meni
    breaklinks=true,        % dozvoljava lomljenje linkova preko redaka
    % bookmarks=true,       % stvara PDF bookmarks (default je true)
    % bookmarksopen=true,   % otvara bookmark panel
    % pdfstartview=FitH,    % kako se PDF otvara
    % pdftitle={Naslov rada},    % naslov PDF-a
    % pdfauthor={Vaše ime},      % autor PDF-a
    % pdfsubject={Tema rada},    % tema PDF-a
}

\begin{document}

\begin{center}
SVEUČILIŠTE JURJA DOBRILE U PULI 

FAKULTET INFORMATIKE

\vspace{45mm} 

\textbf{Tin Pritišanac}

\vspace{20mm} 

\textbf{Strategije iscrtavanja web aplikacija
    kroz različite programske okvire}

\vspace{5mm}
ZAVRŠNI RAD

\vfill

%Upisat tocan mjesec i godinu
Pula, rujan, 2025. godine
\end{center}

\pagenumbering{gobble}
\clearpage
\newpage

\begin{center}
SVEUČILIŠTE JURJA DOBRILE U PULI 

FAKULTET INFORMATIKE

\vspace{45mm} 

\textbf{Tin Pritišanac}

\vspace{20mm} 

\textbf{Strategije iscrtavanja web aplikacija
    kroz različite programske okvire}

\vspace{5mm}
ZAVRŠNI RAD

\end{center}

\vspace{45mm}

\textbf{JMBAG: 0171256219, izvanredni student}

\textbf{Studijski smjer: Informatika}
\bigskip

\textbf{Kolegij: \textbf{Web aplikacije} }

\textbf{Znanstveno područje : Društvene znanosti}

\textbf{Znanstveno polje : Informacijske znanosti}

\textbf{Znanstvena grana : Informacijski sustavi i informatologija}
\bigskip

\textbf{Mentor: doc. dr. sc. Nikola Tanković}

\vfill

\begin{center}

%Upisat tocan mjesec i godinu
Pula, rujan, 2025. godine

\end{center}
\pagenumbering{gobble}
\clearpage
\newpage

\begin{center}
    \begin{figure}[H]
        \centering
        \includegraphics[width=4cm]{slike/uni-logo.jpg}
    \end{figure}

    \vspace{45mm}

    IZJAVA O AKADEMSKOJ ČESTITOSTI

    \vspace{20mm}

    \begin{flushleft}
        Ja, dolje potpisan Tin Pritišanac, ovime izjavljujem da je ovaj seminarski rad rezultat isključivo
        mojega vlastitog rada, da se temelji na mojim istraživanjima te da se oslanja na objavljenu
        literaturu kao što to pokazuju korištene bilješke i bibliografija. Izjavljujem da niti jedan dio
        seminarskog rada nije napisan na nedozvoljen način, odnosno da je prepisan iz kojega
        necitiranog rada, te da ikoji dio rada krši bilo čija autorska prava. Izjavljujem, također, da
        nijedan dio rada nije iskorišten za koji drugi rad pri bilo kojoj drugoj visokoškolskoj,
        znanstvenoj ili radnoj ustanovi.
    \end{flushleft}

    \vspace{20mm}
    \begin{flushright}
        STUDENT~~~~~~~~~~~~~~~ \\
        \bigskip
        \bigskip
        \rule{4cm}{0.4pt}
    \end{flushright}

    \vspace{30mm}

    %Upisat tocan mjesec i godinu
    Pula, rujan, 2025. godine
\end{center}

\pagenumbering{gobble}
\clearpage
\newpage

%Resetirat margine
\restoregeometry

\clearpage
\newpage

\selectlanguage{croatian}
\tableofcontents
\pagenumbering{gobble}
\clearpage
\newpage

\pagenumbering{arabic}

\section{Uvod}

U današnjem digitalno vođenom poslovnom okruženju, sposobnost efikasnog prikupljanja, obrade i analize velikih količina podataka postala je ključni čimbenik uspjeha za organizacije u gotovo svim industrijskim granama. Automobilska industrija, kao jedna od najkompleksnijih i najkonkurentnijih grana gospodarstva, posebno se oslanja na napredne tehnologije skladištenja i analize podataka kako bi razumjela tržišne trendove, potrebe kupaca i operacijske učinkovitosti \cite{Silva2021}.

Skladišta podataka (eng. \textit{data warehouses}) predstavljaju temelj modernih sustava za podršku odlučivanju, omogućujući integraciju različitih izvora podataka u jedinstvenu, koherentnu strukturu optimiziranu za analitičke potrebe \cite{Garani2019}. U kontekstu automobilske industrije, ovakvi sustavi omogućuju analizu složenih odnosa između cijena vozila, karakteristika proizvođača, tržišnih segmenata i regionalnih specifičnosti.

Ovaj rad predstavlja sveobuhvatan pristup razvoju sustava za skladištenje i rudarenje podataka primjenjenog na analizu automobilskog tržišta. Kroz razvoj kompletnog ETL (Extract, Transform, Load) procesa, projekt demonstrira transformaciju sirovih podataka o automobilima u strukturirani dimenzijski model prilagođen OLAP (Online Analytical Processing) analizama. Korištenjem skupa podataka koji sadrži preko 97.000 zapisa o vozilima različitih proizvođača, modela i karakteristika, razvijen je sustav koji omogućuje dubinsku analizu tržišnih trendova i poslovnih uvida.

Glavni cilj ovog projekta je ilustracija praktične primjene teorijskih koncepata skladišta podataka kroz razvoj funcionalnog sustava koji može poslužiti organizacijama poput autoklubova, analitičkih kuća ili samim proizvođačima automobila u donošenju informiranih poslovnih odluka \cite{Nima2018}. Projekt obuhvaća sve ključne faze razvoja sustava - od eksploratorne analize početnih podataka, preko dizajna relacijskog i dimenzijskog modela, do implementacije ETL procesa i prijedloga OLAP analiza.

Struktura rada prati logični tijek razvoja sustava, počevši od analize i pripreme početnog skupa podataka, preko stvaranja normaliziranog relacijskog modela, do konačne implementacije zvjezdastog modela optimiziranog za analitičke potrebe. Svaki korak popraćen je detaljnim objašnjenjima projektnih odluka i praktičnih implementacijskih izazova, čineći ovaj rad korisnim resursom za razumijevanje kompleksnosti razvoja realnih sustava za skladištenje podataka.

Kroz ovaj projekt, nastoji se pokazati kako tehnologije poput Apache Spark-a, MySQL-a i Tableau-a mogu biti integrirane u koherentan sustav koji omogućuje ne samo tehnički ispravan rad, već i stvaranje dodane vrijednosti kroz kvalitetne poslovne uvide \cite{Leka2025}.



\newpage

\section{Projektni zadatak}
\subsection{Cilj i opis projektnog zadatka}
% TODO delete
% \begin{lstlisting}[language=Python, caption={Pisanje Spark DataFrame-a u MySQL skladište podataka}, inputencoding=utf8]
% def drop_fact_tables_first():
%     """Drop fact tables first to avoid foreign key constraint issues"""
%     try:
%         connection = pymysql.connect(
%             host='127.0.0.1',
%             user='root',
%             password='root',
%             database='cars_dw'
%         )
%         cursor = connection.cursor()

%         # Drop fact tables first to remove foreign key constraints
%         fact_tables = ['fact_car_sales']
%         for table in fact_tables:
%             cursor.execute(f"DROP TABLE IF EXISTS {table}")
%             print(f"Dropped fact table {table} (if existed)")

%         connection.commit()
%         cursor.close()
%         connection.close()

%     except Exception as e:
%         print(f"Note: Could not drop fact tables: {e}")
% \end{lstlisting}
\newpage

\section{Odabir i analiza skupa podataka}

\subsection{Izvor podataka}
\subsection{Analiza podataka}
\subsection{Priprema podataka}
\newpage

\section{Relacijski model podataka}

\subsection{Entiteti i atributi}
\subsection{Konceptualni model podataka}
\subsection{Logički model podataka}
\subsection{Kreiranje relacijskog modela podataka}
\subsection{Popunjavanje baze podataka}
\subsection{EER dijagram}
\subsection{Testiranje unosa}
\newpage

\section{Dimenzijski model podataka}

\subsection{Izrada star scheme}
\subsection{Kreiranje dimenzijskih tablica}
\subsection{Kreiranje tablice činjenica}
\newpage

\section{ETL proces}

\subsection{Izvlačenje podataka}
\subsection{Transformacija podataka}
\subsection{Popunjavanje skladišta podataka}
\newpage

\section{OLAP analiza}

\subsection{Definiranje prikaza podataka}
\subsection{Vizualizacija podataka u Tableau}
\subsubsection{Graf 1 TODO}
\newpage

\section{Zaključak}

\newpage

\bibliographystyle{unsrt}
\bibliography{literatura}
\newpage

%Automatski generira listu svih slika
\listoffigures
\newpage

%Automatski generira listu svih tablica
\listoftables
\newpage


\end{document}
